%============================================================================================================
%PARTIE COMPARAISON AVEC DES MODELES AVEC LES DIFFERENT AVION UTILISE 
%================================================================================================================


\chapter{ Comparaison    }\label{ch:structure}

\section{Avions étudiés}

	La distance de décollage dépend fortement des caractéristiques aérodynamiques, de la masse, de la motorisation et des performances globales de l'appareil. Afin d'évaluer l'influence de ces paramètres, plusieurs catégories d'avions ont été sélectionnées pour cette étude.
	
	Le choix des aéronefs s'est porté sur des appareils représentatifs de différents domaines de l'aviation :
	
	\begin{itemize}
		\item le \textbf{Cessna 172}, représentant les avions légers de l'aviation générale ;
		\item l'\textbf{Airbus A320}, représentant les avions commerciaux moyen-courriers ;
		\item le \textbf{Belugar Xl}, représentant les avions de grande capacités
		\item le \textbf{Lockheed Martin C-130J Super Hercules}, représentant les avions militaires de transport ;
		\item le \textbf{Dassault Rafale}, représentant les avions de combat modernes.
	\end{itemize}
	
	Ces appareils présentent des caractéristiques et des performances très différentes, permettant ainsi d'étudier l'influence du poids, de la poussée disponible, de la surface alaire et de la configuration aérodynamique sur la distance de décollage.
	
	Les résultats obtenus pour chaque avion seront ensuite comparés afin d'identifier les principaux paramètres influençant les performances au décollage.
	
	\begin{figure}[H]
		\centering
		
		\begin{subfigure}{0.45\textwidth}
			\centering
			\includegraphics[width=\textwidth]{Cessna172.png}
			\caption{Cessna 172}
		\end{subfigure}
		\hfill
		\begin{subfigure}{0.45\textwidth}
			\centering
			\includegraphics[width=\textwidth]{A320.png}
			\caption{Airbus A320}
		\end{subfigure}
		\vspace{0.5cm}
		
		\begin{subfigure}{0.45\textwidth}
			\centering
			\includegraphics[width=\textwidth]{begura.png}
			\caption{Belugar XL}
		\end{subfigure}
		
		\vspace{0.5cm}
		
		\begin{subfigure}{0.45\textwidth}
			\centering
			\includegraphics[width=\textwidth]{C-130J_super_hercules.png}
			\caption{C-130J Super Hercules}
		\end{subfigure}
		\hfill
		\begin{subfigure}{0.45\textwidth}
			\centering
			\includegraphics[width=\textwidth]{Dassault_rafale.png}
			\caption{Dassault Rafale}
		\end{subfigure}
		
		\caption{Avions retenus pour l'étude comparative des performances au décollage.}
		\label{fig:aircraft_comparison}
	\end{figure}
	
	


	

	
	\subsection{Résultats du Cessna 172}
		
		L'analyse de la distance d'accélération $ Sa $, de la distance de rotation $ Sr $ et de la distance de monté $ Sc $ pour le cas du cessna nous conduire a ces graphes:
		
		\begin{figure}[H]
				
				\begin{subfigure}{0.45\textwidth}
					\centering
					\includegraphics[width=\linewidth]{Sa_cessna.png}
				\end{subfigure}
				\hfill
				\begin{subfigure}{0.45\textwidth}
					\centering
					\includegraphics[width=\linewidth]{Sr_cessna.png}
				\end{subfigure}
				
				\vspace{0.5cm}
				
				\begin{subfigure}{0.45\textwidth}
					\centering
					\includegraphics[width=\linewidth]{Sc_cessna.png}
				\end{subfigure}
				\hfill
				\begin{subfigure}{0.45\textwidth}
					\centering
					\includegraphics[width=\linewidth]{S_cessna.png}
				\end{subfigure}
		\end{figure}
	
	\begin{itemize}
		\item Les méthodes 1, 3 et 4 donnent des distances d'accélération au sol ($S_a$) très proches ; la méthode 2 fournit des valeurs légèrement plus faibles.
		
		\item Les distances de rotation ($S_r$) sont identiques pour toutes les méthodes, celles-ci supposant une rotation de 1 s à la vitesse de décollage.
		
		\item La méthode 2 ne permet pas d'estimer la distance de montée ($S_c$).
		
		\item Les méthodes 3 et 4 produisent des résultats similaires pour $S_c$, tandis que la méthode 1 prédit des distances plus importantes.
		
		\item Les écarts observés résultent principalement des hypothèses retenues pour la modélisation des phases de transition et de montée.
		
		\item De même, les entrées pour la simulation méthode 1 étant différent de celle des autres peut être un cause .
	\end{itemize}
	
	
	
	\subsection{Résultats du Belugar XL}
	
	L'analyse de la distance d'accélération $ Sa $, de la distance de rotation $ Sr $ et de la distance de monté $ Sc $ pour le cas du Belugar XL nous conduire a ces graphes:
	
	\begin{figure}[H]
		
		\begin{subfigure}{0.45\textwidth}
			\centering
			\includegraphics[width=\linewidth]{Sa_belugar.png}
		\end{subfigure}
		\hfill
		\begin{subfigure}{0.45\textwidth}
			\centering
			\includegraphics[width=\linewidth]{Sr_belugar.png}
		\end{subfigure}
		
		\vspace{0.5cm}
		
		\begin{subfigure}{0.45\textwidth}
			\centering
			\includegraphics[width=\linewidth]{Sc_belugar.png}
		\end{subfigure}
		\hfill
		\begin{subfigure}{0.45\textwidth}
			\centering
			\includegraphics[width=\linewidth]{S_belugar.png}
		\end{subfigure}
	\end{figure}
	
	\begin{itemize}
		\item Les quatre méthodes donnent des résultats très proches pour la distance
		 d'accélération au sol ($S_a$), avec des valeurs comprises entre environ 1200 m
		  et 1300 m. Cela montre une bonne cohérence entre les modèles pour cette phase du décollage.
		
		\item Les distances de rotation ($S_r$) sont identiques pour toutes les méthodes ayant fourni un résultat. 
			Cette concordance est due à l'hypothèse commune d'une durée de rotation fixée à 1 seconde à la vitesse de décollage.
		
		\item La méthode 2 ne permet pas d'estimer les distances de montée ($S_c$) et de décollage total ($S$), 
			ce qui limite la comparaison avec les autres modèles.
		
		\item Pour la distance de montée ($S_c$), les méthodes 3 et 4 donnent 
			des résultats très proches, tandis que la méthode 1 prédit une distance
		 	nettement plus élevée. Cette différence traduit une modélisation plus 
		 	conservative de la phase de montée.
		
		\item La distance totale de décollage ($S$) est donc plus importante avec la méthode 1.
		 Les méthodes 3 et 4 conduisent à des estimations plus faibles et très similaires.
		
		\item Globalement, les écarts observés entre les modèles proviennent principalement des
		 hypothèses retenues pour la modélisation de la phase de montée après le décollage,
		  tandis que la phase d'accélération au sol est décrite de manière relativement homogène par l'ensemble des méthodes.
	\end{itemize}
	
	
	\subsection{Commentaires générales }
	
	Les différentes méthodes permettent d'obtenir un ordre de grandeur des composantes de la distance de décollage.
	Bien qu'elles ne fournissent pas une valeur exacte, la comparaison de plusieurs modèles reposant sur des hypothèses
	 différentes permet d'obtenir une bonne estimation plus représentative de la distance réelle de décollage.
	
	